\subsection{A compact ordinary catalog retaining the actual twisted dyadic relations}

\begin{result}{E.187}{Proved}{The exact compact-catalog criterion, and its missing nonemptiness input}
For every rational prime \(p\), fix a nonempty finite set
\[
 \mathcal V_p\subset
 \{(0,0)\}\cup
 \{(e,0),(0,e),(-e,-e):e\in\mathbb Z_{>0}\}.
 \tag{E.187.1}
\]
Assume the following uniform divisibility condition:
for every positive integer \(d\), all entries of every pair
in \(\mathcal V_p\) are divisible by \(d\) for all but
finitely many \(p\). Also choose finitely many allowed real
sign pairs from \((+,+),(+,-),(-,+)\).
These data define a compact catalog of completed Kummer pairs
\[
 A=\epsilon_A\prod_p p^{m_p},\qquad
 B=\epsilon_B\prod_p p^{n_p},\qquad
                     (m_p,n_p)\in\mathcal V_p.
 \tag{E.187.2}
\]
Its everywhere ordinary locus is closed and compact.
One may additionally prescribe any finite set of valuations,
and choose pairwise disjoint finite prime blocks
\(\mathcal B_1,\mathcal B_2,\ldots\), imposing that each
block contain a prime with nonzero valuation pair. These
are closed conditions, and their simultaneous satisfaction
forces infinite support.

There is a compact space parametrizing, up to geometric
conjugation, every full pro-two section whose abelian pair
belongs to this ordinary catalog. It retains the actual
arithmetic Frobenius words, including their path twists.
The existence of such a section satisfying all the block
conditions is equivalent to nonemptiness of every finite
test in one nested sequence of finite tests described below.
Each test uses finitely many catalog choices, one finite
pro-two target, finitely many actual arithmetic relations,
and finitely many ordinary local Kummer tests.

The quantifier is one common catalog followed by every
finite test. Neither one frozen nonzero valuation, nor
separately chosen catalogs at arbitrarily large depths,
supplies this hypothesis. No catalog with all these tests
nonempty is established here. The section conclusion is
at coefficient two; a full pronilpotent conclusion would
also require the other coefficient sections for the same pair.
\end{result}

\begin{proof}
Let \(Y\) be the product of the finite valuation alphabets
and the finite allowed sign set. It is compact. Modulo
\(d\)-th powers, only finitely many factors in E.187.2
can contribute, uniformly over \(Y\). Therefore E.187.2
defines a continuous map from \(Y\) to the completed
global multiplicative group squared. In particular it
does not use compactness of the entire global Kummer group.

For a fixed finite place \(p\), put
\[
 T_p=\{t\in\mathbb Q_p\setminus\{0,1\}:
       (v_p(t),v_p(1-t))\in\mathcal V_p\}.
 \tag{E.187.3}
\]
This is compact: each of its finitely many valuation
patterns bounds \(t\), bounds it away from zero, and
bounds it away from one. Each exact valuation condition
is closed in that bounded region. Its image under the
continuous local Kummer map \(t\mapsto(t,1-t)\) is
therefore compact and closed in the local completed
multiplicative group squared. The ordinary condition at
\(p\) is the inverse image of this closed set. The real
condition is exactly the finite sign condition already
imposed. Intersecting over all \(p\) proves compactness
of the ordinary catalog. A prime-block condition is a
clopen condition on finitely many coordinates of \(Y\).
The blocks being disjoint, satisfying all of them forces
infinitely many nonzero valuations.

We next construct the single global arithmetic presentation
needed for the section parameter space. This construction
is made in the actual group, and does not assume that
separately chosen finite-\(S\) presentations are compatible.
Write
\[
                   G=G_{\mathbb Q}(2),\qquad
                   F=\pi_1(U_{\overline{\mathbb Q}})^{(2)}.
\]
Choose an actual complex conjugation \(c\in G\).
For each odd \(q\), choose a tame inertia generator
\(\tau_q\in G\), normalized by
\(\kappa_2(q)(\tau_q)=1\) in \(\mathbb Z_2\).
The full tame two-primary inertia is present by the
cyclotomic--radical realization in E.185. These elements
tend to one as \(q\) leaves every finite set: every
finite two-extension of \(\mathbb Q\) has only finitely
many ramified primes. This assertion is independent of
the chosen embeddings of the local decomposition groups.

The independent quadratic characters
\(\kappa_2(-1),\kappa_2(2),\kappa_2(q)\) identify the
Frattini quotient of \(G\) with the corresponding compact
product of copies of \(\mathbb F_2\). Surjectivity onto
every finite set of these coordinates follows from
independence of rational prime square classes; compactness
then gives surjectivity onto the product. Choose \(a\in G\)
having value one at \(\kappa_2(2)\) and zero at every
other character in this list. The elements
\[
                         c,a,\tau_q\quad(q\text{ odd})
 \tag{E.187.4}
\]
topologically generate \(G\). Otherwise some finite
two-group quotient would have their images in a proper
subgroup, and then in a subgroup of index two; a nonzero
quadratic character would vanish on every element in
E.187.4, contradicting the dual basis just constructed.

Let \(\mathcal F\) be the free pro-two group on the
converging pointed family \(C,A,X_q\), with
\(X_q\to1\). Equivalently its finite quotients are
specified by arbitrary images of \(C,A\) and by images
of the \(X_q\) that are one for all but finitely many
\(q\). Its universal property concerns exactly converging
families of images. Sending the generators to E.187.4
gives a continuous surjection \(\mathcal F\to G\).
For each \(q\), choose an actual arithmetic Frobenius
\(\sigma_q\) in its chosen local decomposition group
and any preimage \(Y_q\in\mathcal F\). Then a defining
family of closed normal relations is
\[
       C^2,\qquad R_q=Y_qX_qY_q^{-1}X_q^{-q}
                                    \quad(q\text{ odd}).
 \tag{E.187.5}
\]

Here is a justification of the infinite presentation.
Finite embedding problems for \(\mathcal F\) lift
generator by generator: only finitely many generators
have nontrivial images in a finite quotient, and the
remaining generators may be lifted to one. Hence
\(H^2(\mathcal F,\mathbb F_2)=0\).
Continuous finite-coefficient cohomology commutes with
the filtered union of the finite-ramification quotients:
a cocycle, and any coboundary witnessing an equality,
factor through a finite quotient, whose field has finite
ramification. Taking the filtered union of E.185.1 gives
\[
 H^2(G,\mathbb F_2)
       \simeq\mathbb F_2\oplus
                 \bigoplus_{q\text{ odd}}\mathbb F_2,
 \tag{E.187.6}
\]
where the coordinates are the real and odd local invariants
and the omitted dyadic invariant is fixed by reciprocity.
The generators E.187.4 give an isomorphism on first
cohomology. Thus the free-presentation five-term sequence
and the local relator evaluations identify E.187.6 with
the continuous dual of the relation module. Evaluating
on the words E.187.5 gives precisely its displayed
coordinates. Their closed span is consequently the
whole relation module. The finite-two-group normal-
generation argument in E.185 now shows that their closed
normal closure is the full kernel. It applies unchanged
to an infinite family, since any proposed nontrivial
remaining normal subgroup is detected in a finite
two-group quotient. Finally \(R_q\to1\): in every
finite quotient, \(X_q=1\) for all sufficiently large
\(q\), independently of the values of \(Y_q\).
This proves E.187.5 with its required continuity.

Let \(\rho:G\to\operatorname{Aut}(F)\) be the actual
tangential action. It is unramified at every odd prime,
so \(\rho(\tau_q)=1\). Its image is pro-two. Indeed
the action on \(F/\Phi(F)\) is the identity, and the
kernel of the automorphism action on that quotient is
pro-two: on each finite lower exponent-two quotient,
its successive congruence kernels embed in groups of
homomorphisms between elementary abelian two-layers.
Thus any full pro-two section cocycle on \(G_{\mathbb Q}\)
factors through \(G\), by applying the maximal-pro-two
universal property to its graph in
\(F\rtimes\rho(G)\). The actual unramified action used
here is the based action already verified in A.31.

Fix positive peripheral generators \(x,y\) of \(F\),
and put \(z=(xy)^{-1}\) for the third peripheral class.
To a valuation pair in \(\mathcal V_q\) assign respectively
the following element:
\[
 (e,0)\mapsto x^e,\quad
 (0,e)\mapsto y^e,\quad
 (-e,-e)\mapsto z^e,\quad
 (0,0)\mapsto1.
 \tag{E.187.7}
\]
Let \(v_q\) denote this element. Consider the compact
parameter space
\[
             \mathcal P=Y\times F^2\times\prod_{q\text{ odd}}F.
 \tag{E.187.8}
\]
Write its last coordinates as \(u_c,u_a,g_q\). Assign
to the free generators of \(\mathcal F\) the graph values
\[
 C\mapsto(u_c,\rho(c)),\qquad
 A\mapsto(u_a,\rho(a)),\qquad
 X_q\mapsto(g_qv_qg_q^{-1},1).
 \tag{E.187.9}
\]
Uniform divisibility makes these last values tend to one
uniformly over the entire parameter space. In any finite
two-group quotient of \(F\), a sufficiently large power
of two annihilates every element; all the exponents in
all but finitely many alphabets are divisible by it.
Conjugating cannot change this conclusion. Thus E.187.9
always defines a continuous homomorphism from
\(\mathcal F\) to \(F\rtimes\rho(G)\).

Moreover evaluation of every fixed word of \(\mathcal F\)
is continuous as a function on \(\mathcal P\). After
passing to any finite characteristic quotient of \(F\)
and its finite arithmetic action, only finitely many
\(X_q\)-images can be nontrivial, uniformly in the
parameters. Evaluation there is a function of finitely
many finite coordinates. In particular imposing all
relations E.187.5 is a closed condition. Their arithmetic
coordinates already vanish, so they assert exactly that
the graph descends to an actual cocycle \(G\to F\).

It remains to impose the chosen abelian pair, rather
than just its inertia valuations. Choose compatible
Kummer cocycles for the positive rational primes and
for minus one. The series
\[
 \kappa_2(A)=\kappa_2(\epsilon_A)+\sum_p m_p\kappa_2(p),
 \qquad
 \kappa_2(B)=\kappa_2(\epsilon_B)+\sum_p n_p\kappa_2(p)
 \tag{E.187.10}
\]
converge uniformly on \(G\), uniformly over \(Y\).
Their values lie in the fixed lattice \(\mathbb Z_2(1)\),
and the scalar coefficients tend uniformly to zero.
They are therefore continuous cocycles depending
continuously on the catalog parameters. At \(\tau_q\)
their values are exactly \((m_q,n_q)\): other prime
Kummer cocycles are unramified there, and the chosen
tame generator has the normalization stated above.
Impose equality of the abelianizations of \(u_c,u_a\)
with E.187.10 at \(c,a\). These are closed conditions,
and equality on the generators then gives literal
equality of the entire abelian cocycles.

Every section with this abelian cohomology class admits
such a representative: the difference from the specified
abelian cocycle is a coboundary in \(\mathbb Z_2(1)^2\),
and its conjugating element lifts to \(F\).
For an ordinary pair, at each nonzero odd support prime
A.31 identifies its full local class with the ordinary
boundary point class. Its tame image is consequently a
conjugate of E.187.7. At a zero-valuation odd prime its
inertia image is trivial even without asserting full
local point matching: if an inertia element had first
nonzero lower-central weight \(w\geq2\), the tame
relation would multiply that nonzero torsion-free
coordinate by both \(q\) and \(q^w\), a contradiction.
Wild inertia has trivial pro-two image. Thus every
ordinary full pro-two section is represented in the
closed locus of \(\mathcal P\) just constructed.
Conversely every parameter in that locus is an actual
section with the specified abelian pair. Adding the
ordinary and block conditions proves the compactness
assertion with all local frames retained.

We finish by making the finite-test quantifier precise.
Use the lower exponent-two series of \(F\), and let
\(Q_j\) be its finite characteristic quotient at stage
\(j\). At test \(j\), require the following conditions:
the first \(j\) block conditions; the abelian equalities
at \(c,a\) modulo \(2^j\); all graph relations in
E.187.5 modulo \(Q_j\); and, at the first \(j\)
rational primes, ordinary local Kummer membership
modulo \(j!\), using the compact sets \(T_p\).
Choose the quotients and enumerate the primes so that
these tests are nested.

Each test depends on only finitely many finite coordinates.
If the exponent of \(Q_j\) is \(2^{e_j}\), then all
relations \(R_q\) with every pair in \(\mathcal V_q\)
divisible by \(2^{e_j}\) hold automatically: their
\(X_q\)-image is one. There are only finitely many
remaining \(q\). The ordinary tests modulo \(j!\)
also depend on only finitely many catalog coordinates
by uniform divisibility. Their local images are finite.
Hence each test defines a clopen compact subset of
\(\mathcal P\). If all the test sets are nonempty,
their nested intersection is nonempty. The intersection
satisfies every group relation and abelian equality.
At each local prime, compactness of \(T_p\) makes
membership in all its finite Kummer images equivalent
to membership in its full Kummer image. Thus the
intersection gives the asserted everywhere ordinary
full pro-two section, with infinite support when the
block conditions were included. The converse follows
by projecting any such full section to the finite tests.

Fixing only one valuation does not furnish a substitute
for the uniform catalog. For example, for varying odd
primes \(q\ne23\), the rational points
\(t_q=23q\) give ordinary pairs and full point sections
with fixed valuation pair \((1,0)\) at twenty-three.
Their first global square classes \([23q]\) are all
distinct. The finite-coefficient global Kummer group
has its discrete continuous-cochain topology, so this
family cannot lie in a compact space with continuous
abelian projection. This example already consists of
full point sections.

Finally E.103 shows that the rational Kummer locus is
closed and discrete: a fixed fourth-power fiber lies
on a smooth genus-three Fermat twist and is finite.
Its intersection with any compact ordinary catalog
is therefore finite. One may delete clopen neighborhoods
of these finitely many rational pairs to obtain a compact
ghost sublocus when one remains. That operation need
not preserve nonemptiness of the finite section tests.
Singleton alphabets recover fixed-pair compactness;
the genuinely missing input is nonemptiness at every
depth for one catalog admitting infinite support.
\end{proof}
